% Fragmento: requiere amsmath, booktabs y bibliografía de esta revisión.
\section{Verificación resistente condicionada del eje transversal M8}
\label{sec:M8_Kt_revision}
Se estudia exclusivamente la varilla M8$\times$1,25, con clase resistente
\textbf{8.8 candidata} y capa de biomasa de \textbf{4 mm}. Se conservan el
historial acoplado 3C de 1440 estados, la presión de Holm de
2,25867 MPa, el coeficiente efectivo de torque de 0,1 y
$E=210$ GPa. El doble empotramiento continúa siendo una hipótesis del
modelo, pendiente de demostrar mediante el diseño de las fijaciones en ABS.
No se modificaron el CAD ni los resultados históricos.

La sección circular equivalente del modelo emplea el diámetro de núcleo
de referencia $d_r=6,466$ mm, con $I_r=\pi d_r^4/64$ y
$\sigma_{b,\mathrm{nom},r}=32\sqrt{M_y^2+M_z^2}/(\pi d_r^3)$.
Los resultados heredados son:
\begin{center}
\begin{tabular}{lrr}
\toprule
Magnitud máxima independiente & Valor & Estado del historial\\
\midrule
Fuerza total de compactación & 918,00 N & 476\\
Torque resistente del accionamiento & 2,7804 N\,m & 476\\
Flecha de la varilla & 0,1569 mm & 301\\
Tensión nominal de flexión & 315,36 MPa & 541\\
\bottomrule
\end{tabular}
\end{center}
No se combinan estos máximos como si fueran simultáneos. La sección de
máxima tensión está en el empotramiento derecho, a 128,1178 mm del extremo
de referencia de la varilla, para una posición angular de 135,25 grados.
El momento flector resultante correspondiente es 8,3698 N\,m.
El torque de accionamiento alrededor del eje de la matriz no se añade
como torsión alrededor del eje transversal.

\subsection{Concentración de tensiones y alcance de la comprobación}
La auditoría de EP2 identifica los valores $K_t=1,5$ y $2,5$ como
sensibilidades paramétricas, sin selección mediante un radio de raíz,
tabla o ecuación correspondiente a la rosca M8\cite{ep2_auditoria_Kt}.
No existe información suficiente para seleccionar un $K_t$ único de esta
varilla y de su fijación. La geometría local de raíz y el contacto de los
filetes son relevantes\cite{pedersen_2013_roscas}; no están definidos
en el STEP ni en las especificaciones del producto todavía no adquirido.

Se define el factor que necesitaría esta memoria como
\begin{equation}
 K_{t,r}=\frac{\sigma_{b,\mathrm{local},\mathrm{el}}}
 {\sigma_{b,\mathrm{nom},r}}.
 \label{eq:Kt_raiz_M8}
\end{equation}
El subíndice $r$ identifica una \emph{tensión nominal calculada con el
núcleo}; no representa un valor tabulado. Así se evita volver a incluir
la reducción de diámetro ya incorporada en los 315,36 MPa.
La clase 8.8 aporta condicionalmente $R_{p0,2,\min}=640$ MPa para el
diámetro considerado, si se confirma la conformidad del producto
\cite{boellhoff_manual_8100_2023}. Para una aproximación local de flexión
uniaxial elástica, resulta
\begin{equation}
 n_{y,\mathrm{flex}}(K_{t,r})
 =\frac{640}{315,361913\,K_{t,r}}
 =\frac{2,029414}{K_{t,r}}.
 \label{eq:ny_M8_condicionado}
\end{equation}
El valor 2,0294 es el cociente \emph{nominal sin concentración}, no un
factor de seguridad aprobado de la varilla roscada. El umbral
$K_{t,r}=2,0294$ corresponde a $n_{y,\mathrm{flex}}=1$ y al inicio
teórico de fluencia según esta aproximación; tampoco constituye una
selección de $K_t$ ni un criterio de seguridad de diseño acordado.
Por encima del umbral, una extrapolación elástica no describe la
redistribución plástica ni permite afirmar colapso.

\subsection{Conclusión de diseño}
La alternativa M8 se conserva como preferente, con clase 8.8 candidata,
\textbf{sin aprobación resistente}. Se requieren el perfil y radio real
de raíz, tolerancias del producto, ubicación de los filetes respecto de
las fijaciones, geometría de engrane y precarga para justificar el
concentrador en la sección crítica. La equivalencia de viga con núcleo
sirve como referencia nominal global; no resuelve los esfuerzos de una
rosca helicoidal ni demuestra el empotramiento real en ABS.

El factor efectivo de fatiga $K_f$ permanece sin seleccionar. Para cerrar
fatiga faltan además fabricación y acabado de la rosca, propiedades
cíclicas representativas, sensibilidad a la entalla y espectro/vida de
servicio. La comprobación deberá utilizar tensiones firmadas en una misma
fibra, conservando la correlación temporal; el eje fijo no implica
flexión rotativa completamente alternante. El anexo desarrolla las
definiciones, la auditoría de fuentes y las sensibilidades sin asignar
un intervalo físico no demostrado a $K_t$.
